\subsection{Log Rules to Separate} 
One use for the log rules is to rewrite the log of a complex expression as a combination of logs of much simpler expressions.
\begin{note}
In the log rules, we assume that every expression whose log we take is positive, since the log is only defined for positive numbers. Sometimes, we will include an absolute value to make this less of a concern. In our work with log rules in this class, we will assume the values of all variables are such that there is no concern about domain.
\end{note}
\begin{note}
We have no log rules to split up \makeblank[1in]. If we have \makeblank[1in] (or \makeblank[1in]) inside a log, we have to leave it alone.
\end{note}
\begin{example}
Rewrite $\log_7\left(\dfrac{2x^4\sqrt{x-4}}{y^5z}\right)$ using the simplest possible logs.
\begin{lpic}{}{9}
\foreach \y in {-3,-5,-7,-9} \regtextleft{0}{\y}{${}={}$};
\end{lpic}
\end{example}
We notice the following, which gives us a shortcut.
\begin{itemize}
\item Factors in the \makeblank[1.5in] give us positive terms
\item Factors in the \makeblank[1.5in] give negative terms
\end{itemize}
Let's see another example, using that shortcut.
\begin{example}
Rewrite $\log_4\left(\dfrac{3(x-1)^5\sqrt[3]{y}}{z^4\sqrt{y+1}}\right)$ using the simplest possible logs.
\begin{lpic}{}{3}
\foreach \y in {-1,-3} \regtextleft{0}{\y}{${}={}$};
\end{lpic}
\end{example}